\documentclass[11pt]{article}
\usepackage[T1]{fontenc}
\usepackage[utf8]{inputenc}
\usepackage{lmodern}
\usepackage{amsmath,amssymb}
\usepackage[margin=1in]{geometry}
\usepackage{booktabs}
\usepackage{xcolor}
\usepackage{pgfplots,pgfplotstable}
\usepgfplotslibrary{groupplots}
\usepackage[hidelinks]{hyperref}
\pgfplotsset{compat=1.18}
\definecolor{baseline}{HTML}{24557A}
\definecolor{splitcolor}{HTML}{A3442B}
\setlength{\parskip}{0.35em}
\setlength{\parindent}{0pt}
\emergencystretch=2em
\newcommand{\E}{\mathbb E}
\newcommand{\R}{\mathbb R}
\newcommand{\KL}{\operatorname{KL}}
\newcommand{\Var}{\operatorname{Var}}
\newcommand{\D}{\mathrm d}
\pgfplotsset{membership axis/.style={
 width=0.44\textwidth,height=0.30\textwidth,
 xmin=0.65,xmax=4.35,ymin=0.40,ymax=1.045,
 xtick={1,2,3,4},xticklabels={0,0.1,1,4},ytick={0.4,0.6,0.8,1.0},
 grid=major,grid style={gray!16},tick label style={font=\small},
 label style={font=\small},title style={font=\small},
 legend style={font=\scriptsize,draw=none,fill=white},legend cell align=left}}
\pgfplotstableread{
@@REGULAR64@@
}\regularsmall
\pgfplotstableread{
@@REGULAR256@@
}\regularlarge
\pgfplotstableread{
@@ALTERNATING64@@
}\alternatingsmall
\pgfplotstableread{
@@ALTERNATING256@@
}\alternatinglarge
\newcommand{\membershipplot}[1]{%
\addplot+[color=baseline,mark=*,mark size=1.5pt,error bars/.cd,y dir=both,y explicit]
 table[x expr=\thisrow{index}-0.06,y=lr,y error minus=lrminus,y error plus=lrplus]{#1};
\addplot+[color=splitcolor,mark=diamond*,mark size=1.8pt,error bars/.cd,y dir=both,y explicit]
 table[x expr=\thisrow{index}+0.06,y=split,y error minus=splitminus,y error plus=splitplus]{#1};
\addplot[gray,dashed,forget plot] coordinates {(0.65,0.95) (4.35,0.95)};
}

\title{Inference for Weak OU Mean Reversion:\\
A Reproducible Finite-Horizon Study}
\author{Junguang Jia}
\date{}

\begin{document}
\maketitle

\section*{Abstract}
A short trajectory can provide little information about weak mean reversion,
even when observations are frequent. This study examines that difficulty for a
scalar Ornstein--Uhlenbeck process with unknown drift and diffusion scale,
conditional on a fixed initial state. It derives the exact transition likelihood,
profiles the nuisance parameters analytically, and implements a chronological
split-profile confidence rule based on universal inference. A restricted
continuous-path Kullback--Leibler calculation distinguishes an information limit
from the behavior of a particular procedure. In a saved development experiment
with 1,600 trajectories, the split rule has true-parameter membership between
0.99 and 1.00 but excludes zero in at most 1\% of each cell. An uncalibrated
chi-square likelihood-ratio diagnostic has membership between 0.57 and 0.91.
The Julia implementation is checked against frozen Python and separately coded
base-R references. The result is a reproducible technical and numerical study;
it does not introduce a new inference method or establish certified confidence
endpoints or a calibrated bootstrap comparison.

\section{Introduction}
For a persistent continuous-time process, a positive fitted mean-reversion
coefficient does not by itself demonstrate that reversion is well identified.
The relevant inferential questions include whether zero can be excluded,
whether arbitrarily slow reversion remains plausible, and whether uncertainty
has a finite upper endpoint. These questions become especially delicate when
both the drift intercept and diffusion scale must be estimated from one finite
observation window.

This report studies
\begin{equation}
 \D X_t=(a-\kappa X_t)\,\D t+\sigma\,\D W_t,
 \qquad a\in\R,\quad\kappa\geq0,\quad\sigma>0,
 \label{eq:model}
\end{equation}
conditional on an observed, fixed $X_0=x_0$. Observation times are deterministic,
or exogenous so that conditioning on the schedule preserves the stated
transition law. At $\kappa=0$ the model is Brownian motion with drift $a$.
For $\kappa>0$, the long-run mean is $\theta=a/\kappa$; it is not identified
at zero. A bounded-$\theta$ restriction would force $a\to0$ as $\kappa\to0$
and would therefore change the boundary experiment.

Diffusion drift bias and persistent-process inference have substantial prior
literature~\cite{tangchen,yu}. Hansen's grid bootstrap and its continuous-time
extension by Lui, Xiao and Yu supply important comparison
methods~\cite{hansen,lui}. The split-profile construction below is a
specialization of universal inference~\cite{universal}. The present objective
is to verify a boundary-aware implementation and expose limitations that a
subsequent calibrated comparison must address. Migration to Julia is an
implementation choice, not a methodological contribution.

\section{Exact likelihood and inferential methods}
\subsection{Transition law and nuisance profiling}
Let $0=t_0<t_1<\cdots<t_n=T$ and $\Delta_i=t_i-t_{i-1}>0$. Solving the
linear stochastic differential equation gives
\begin{equation}
 X_i\mid X_{i-1}\sim
 N\!\left(\phi_iX_{i-1}+ab_i,\;\sigma^2v_i\right),
 \quad
 \phi_i=e^{-\kappa\Delta_i},\quad
 b_i=\frac{1-e^{-\kappa\Delta_i}}{\kappa},\quad
 v_i=\frac{1-e^{-2\kappa\Delta_i}}{2\kappa}.
 \label{eq:transition}
\end{equation}
At zero, the continuous limits are exactly
$\phi_i=1$ and $b_i=v_i=\Delta_i$. The implementation uses stable exponential
expressions and an explicit zero branch, with no Euler approximation or
positive lower bound on $\kappa$. There is no stationary initial density.

For a candidate $\kappa$, write $r_i=x_i-\phi_i x_{i-1}$ and $s^2=\sigma^2$.
The conditional log likelihood is
\[
 \ell(\kappa,a,s^2)=-\frac12\sum_{i=1}^n
 \left\{\log(2\pi s^2v_i)+\frac{(r_i-ab_i)^2}{s^2v_i}\right\}.
\]
Weighted least squares and one-dimensional variance maximization yield
\begin{align}
 \widehat a_\kappa
 &=\frac{\sum_i b_ir_i/v_i}{\sum_i b_i^2/v_i},&
 R_\kappa&=\sum_i\frac{(r_i-\widehat a_\kappa b_i)^2}{v_i},
 \label{eq:profile-estimates}\\
 \widehat\sigma^2_\kappa&=\frac{R_\kappa}{n},&
 \ell_p(\kappa)&=-\frac12\left[
 n\{\log(2\pi\widehat\sigma^2_\kappa)+1\}+\sum_i\log v_i\right].
 \label{eq:profile}
\end{align}
Indeed, the residual quadratic is
$R_\kappa+(a-\widehat a_\kappa)^2\sum_i b_i^2/v_i$, with a strictly positive
coefficient. When $R_\kappa>0$, differentiating the remaining likelihood in
$s^2$ gives the maximum at $R_\kappa/n$. The divisor counts transitions; it is
not the residual degrees-of-freedom divisor of an unbiased variance estimator.
When $R_\kappa=0$, the likelihood supremum is infinite as $s^2\downarrow0$.
This degeneracy is retained rather than replaced by a variance floor.

The numerical fit evaluates zero and a positive logarithmic grid up to $200/T$, then refines
local candidates by a bounded scalar search. This is not a certified global
maximum. As $\kappa\to\infty$, the profiled likelihood tends to the free-mean,
free-variance iid Gaussian likelihood of $x_1,\ldots,x_n$. To see this,
reparameterize $a=\kappa\beta$ and $\sigma^2=2\kappa\tau^2$ in
\eqref{eq:transition}; the means tend to $\beta$ and variances to $\tau^2$.
The explicit weighted least-squares profiles converge to the successor sample
mean and variance, giving convergence of the profiled likelihood as well.
If the observed successor states are all equal, the limiting supremum is
infinite. Cap contact and comparison with this tail limit are reported as
warnings. A search cap is never interpreted as a confidence endpoint.

\subsection{An explicitly uncalibrated likelihood-ratio diagnostic}
Let $\widetilde\ell_{\max}$ denote the best value returned by the bounded
search. The diagnostic statistic is
\[
 D(\kappa)=2\{\widetilde\ell_{\max}-\ell_p(\kappa)\}.
\]
The development comparison accepts a candidate when
$D(\kappa)\leq\chi^2_{1,0.95}=3.84145882$. This cutoff is retained to examine
its behavior; it is not justified here by Wilks' theorem. Neither an ordinary
chi-square limit nor a simple boundary mixture is assumed reliable in the
fixed-span, weak-reversion regime. Approximate global optimization is an
additional numerical limitation of this diagnostic.

\subsection{Chronological split-profile rule}
Choose a fixed split after $m$ transitions. Fit an OU density on the prefix
$X_0,\ldots,X_m$, freeze its parameters, and evaluate the remaining transition
product $q(X_{m+1:n}\mid\mathcal F_m)$, where
$\mathcal F_m=\sigma(X_0,\ldots,X_m)$ with the schedule conditioned on.
For each realized prefix, $q$ must be a normalized conditional density. Define
\begin{equation}
 E_\kappa=
 \frac{q(X_{m+1:n}\mid\mathcal F_m)}
 {\sup_{a\in\R,\,\sigma>0}
 p_{\kappa,a,\sigma}(X_{m+1:n}\mid X_m)},
 \qquad
 C_\alpha=\{\kappa\geq0:E_\kappa\leq1/\alpha\}.
 \label{eq:evalue}
\end{equation}
The denominator is the analytic nuisance supremum on the suffix. At the true
parameters $(\kappa_0,a_0,\sigma_0)$, conditional-density domination gives
\[
 \E_{\kappa_0,a_0,\sigma_0}[E_{\kappa_0}\mid\mathcal F_m]
 \leq\int q(y\mid\mathcal F_m)\,\D y=1.
\]
Markov's inequality therefore implies
$P_{\kappa_0,a_0,\sigma_0}(\kappa_0\in C_\alpha)\geq1-\alpha$.
Conditioning on $X_m$ handles prefix/suffix dependence; independence of the two
segments is not required. Approximate prefix fitting preserves this argument
when it still produces a normalized positive-variance density. An underestimated
denominator, future-informed fitting, or an unreported training failure does
not inherit the guarantee.

This is the conditional/profile-likelihood principle of universal
inference~\cite{universal}. Its mathematical guarantee concerns the exact model
and exact nuisance supremum. Floating-point agreement is not a certificate that
every computed denominator is a conservative upper bound. The rule is
fixed-horizon and model-based, not a confidence sequence, prediction interval,
or distribution-free guarantee.

\section{Information limitations over a finite horizon}
\subsection{Dimensionless parameters}
Brownian scaling of $U_s=(X_{Ts}-x_0)/(\sigma\sqrt T)$ gives
\[
 \D U_s=(\delta-cU_s)\,\D s+\D B_s,
 \qquad c=\kappa T,\qquad
 \delta=\frac{(a-\kappa x_0)\sqrt T}{\sigma}.
\]
Thus $c$ measures reversion relative to the observation span, but does not
by itself determine every experiment: $\delta$ and the normalized schedule
also matter. Increasing $n$ at fixed $T$ is an infill experiment; increasing
$T$ at a fixed observation step is a different information change.

\subsection{A restricted pathwise KL calculation}
For this subsection only, assume $a=0$, $x_0=0$, and the same known $\sigma$
under both laws. Let $P_\kappa$ and $P_0$ denote continuous-path laws on $[0,T]$
for OU and driftless Brownian motion. Under $P_0$, $X_t=\sigma W_t$. It\^o's
formula expresses the candidate likelihood ratio as
\[
 L_t=\exp\left\{-\frac\kappa\sigma\int_0^tX_s\,\D W_s
 -\frac{\kappa^2}{2\sigma^2}\int_0^tX_s^2\,\D s\right\}
 =\exp\left\{\frac{\kappa t}{2}-\frac{\kappa X_t^2}{2\sigma^2}
 -\frac{\kappa^2}{2\sigma^2}\int_0^tX_s^2\,\D s\right\}.
\]
For $\kappa\geq0$, this positive local martingale is bounded by
$e^{\kappa T/2}$ on the finite horizon. Localization and dominated convergence
make it a true martingale. Girsanov's theorem then identifies
$L_T=\D P_\kappa/\D P_0$. Under $P_\kappa$, write
\[
 \log L_T=-\frac\kappa\sigma\int_0^T X_t\,\D W_t^{(\kappa)}
 +\frac{\kappa^2}{2\sigma^2}\int_0^T X_t^2\,\D t,
\]
where $W^{(\kappa)}$ is Brownian motion under $P_\kappa$. Its stochastic
integral has finite second moment and hence mean zero under this measure. Using
$\E_\kappa X_t^2=\sigma^2(1-e^{-2\kappa t})/(2\kappa)$ gives
\begin{align}
 \KL(P_\kappa\Vert P_0)
 &=\frac{\kappa^2}{2\sigma^2}\int_0^T\E_\kappa X_t^2\,\D t\\
 &=\frac c4-\frac{1-e^{-2c}}8
 =\frac{c^2}{4}-\frac{c^3}{6}+O(c^4).
 \label{eq:kl}
\end{align}
The identity is continuous at $c=0$.

For discrete observations, the Gaussian chain rule gives
\begin{equation}
 \KL(P_\kappa^{\rm obs}\Vert P_0^{\rm obs})
 =\frac12\sum_i\left\{
 \frac{v_i}{\Delta_i}-1-\log\frac{v_i}{\Delta_i}
 +\frac{(\phi_i-1)^2 V_{i-1}}{\Delta_i}\right\},
 \quad
 V_{i-1}=\frac{1-e^{-2\kappa t_{i-1}}}{2\kappa}.
 \label{eq:sample-kl}
\end{equation}
At $\kappa=0$, take $V_{i-1}=t_{i-1}$ by continuity.
Sampling is a measurable function of the full path. Data processing therefore
bounds sampled KL by path KL and implies monotonicity for nested refinement.
Any test of the restricted Brownian null with size at most $\alpha$ has power
bounded by
\begin{equation}
 \operatorname{Power}\leq
 \min\left\{1,\alpha+\sqrt{\KL(P_\kappa\Vert P_0)/2}\right\}.
 \label{eq:pinsker}
\end{equation}
This follows from total variation and Pinsker's inequality. At $c=0.1$,
path KL is approximately $0.0023413441$, so the level-$0.05$ power bound is
$0.0842151$. Arbitrarily dense sampling cannot exceed the full-path information
in this restricted experiment. The bound is loose and need not be attainable;
it is not a power curve for the unknown-nuisance split rule. A difference between
that rule's power and this envelope does not identify the cost of splitting.

\begin{table}[!htbp]
\centering\small
\caption{Restricted information calculation with $a=x_0=0$ and known $\sigma$.
Discrete KL uses eight equal gaps; the final column is the continuous-path
level-$0.05$ power upper bound. These are deterministic values, not simulations.}
\label{tab:information}
\begin{tabular}{rrrr}
\toprule
$c$ & Sampled KL ($n=8$) & Path KL & Power upper bound\\
\midrule
@@INFORMATION_TABLE@@
\bottomrule
\end{tabular}
\end{table}

\section{Development experimental design}
The main empirical results use a saved Julia development pilot with
$T=1$, $a=x_0=0$, $\sigma=1$, $\kappa\in\{0,0.1,1,4\}$, and
$n\in\{64,256\}$ transitions. Regular gaps are $1/n$. The deterministic
alternating schedule repeats $0.25/n$ and $1.75/n$, preserving the same $T$
and $n$. This gives 16 cells, each with $R=100$ independently simulated
trajectories, for 1,600 trajectories in total. Both $a$ and $\sigma$ are
unknown to the inference procedures, even though their generating values are
fixed. The split is $m=n/2$ and $\alpha=0.05$.

Exact transition simulation uses a fresh Julia Xoshiro generator for each
trajectory. The base seed is $2026092803$; the stored seed equals this value
plus the trajectory index in the declared loop order. Different resolutions
and schedules use different paths. These comparisons are therefore not paired
nested-path infill experiments. A future infill design must derive coarse
observations from a common exact fine trajectory; longer-span comparisons
should use common path prefixes.

For each trajectory, the pilot evaluates true-parameter membership and exclusion
of zero directly from the two statistics. It does not construct full confidence
sets or measure their widths. With $I_r$ the relevant binary outcome,
$\widehat p=R^{-1}\sum_r I_r$ and the reported plug-in Monte Carlo standard
error is $\sqrt{\widehat p(1-\widehat p)/R}$. The uncertainty unit is one
independent trajectory, not an observation within a trajectory. Pointwise 95\%
Wilson intervals are also reported, including at zero or one observed rates;
these describe simulation uncertainty, not parameter confidence sets for $\kappa$.

All attempted trajectories remain in denominators. If $S$ observed successes and
$F$ unresolved trajectories occur in $R$ attempts, the empirical success fraction
is bounded by $[S/R,(S+F)/R]$, with failures reported separately. The saved Julia
pilot has zero failed trajectories and no recorded cap, tail, optimizer or
numerical warnings. This does not certify global optimization. The entire pilot
is exploratory development evidence, with no locked confirmation design or
method selection justified by these results.

\section{Results and numerical verification}
\subsection{Coverage alone does not establish informativeness}
Table~\ref{tab:results} reports the regular $n=256$ cells; Figure~\ref{fig:membership}
shows true-parameter membership in all 16 cells. The chi-square LR diagnostic
has membership between $0.57$ and $0.91$, whereas split-profile membership is
between $0.99$ and $1.00$. At the four $\kappa=0$ cells, the LR diagnostic falsely
excludes zero in 28\%, 43\%, 30\% and 34\% of trajectories (regular $n=64,256$,
then alternating $n=64,256$). This is strong evidence of poor calibration for
that particular cutoff in the tested development settings; it does not show
that all likelihood-based inference fails.

\begin{table}[!htbp]
\centering\small
\caption{Selected Julia development results: regular schedule, $n=256$ and
100 trajectories per row. Membership is at the true $\kappa$; zero exclusion
is a false rejection at $\kappa=0$ and empirical power otherwise. LR uses an
uncalibrated cutoff. Fitted $\kappa$ values use bounded numerical search.}
\label{tab:results}
\begin{tabular}{rrrrrr}
\toprule
 & & \multicolumn{2}{c}{Truth membership} & \multicolumn{2}{c}{Zero exclusion}\\
\cmidrule(lr){3-4}\cmidrule(lr){5-6}
$\kappa$ & Mean $\widehat\kappa$ & LR & Split & LR & Split\\
\midrule
@@RESULTS_TABLE@@
\bottomrule
\end{tabular}
\end{table}

\begin{figure}[!ht]
\centering
\begin{tikzpicture}
\begin{groupplot}[membership axis,group style={group size=2 by 2,horizontal sep=0.8cm,vertical sep=1.1cm}]
\nextgroupplot[title={Regular, $n=64$},ylabel={Truth membership},legend pos=south east]
\membershipplot{\regularsmall}
\legend{LR diagnostic,Split profile}
\nextgroupplot[title={Regular, $n=256$}]
\membershipplot{\regularlarge}
\nextgroupplot[title={Alternating, $n=64$},xlabel={True $\kappa T$},ylabel={Truth membership}]
\membershipplot{\alternatingsmall}
\nextgroupplot[title={Alternating, $n=256$},xlabel={True $\kappa T$}]
\membershipplot{\alternatinglarge}
\end{groupplot}
\end{tikzpicture}
\caption{True-parameter membership in the Julia development pilot. Bars are
pointwise 95\% Wilson intervals from 100 independent trajectories per cell;
the dashed line is 0.95. Horizontal positions are categorical, slightly offset
for legibility; connecting segments guide the eye and do not interpolate power.
Blue circles show the uncalibrated LR diagnostic; rust diamonds show split profile.}
\label{fig:membership}
\end{figure}

The split rule excludes zero in at most one of the 100 trajectories in any cell,
including at $c=4$. No exclusions in a cell give the Wilson interval
$[0,0.03699]$; 100 memberships give $[0.96301,1]$. Neither zero estimated MCSE
at these extremes nor very high membership certifies exact empirical risk.
The inability to exclude zero is a substantive negative informativeness result.
It motivates separating intrinsic weak information, nuisance estimation and
chronological splitting, rather than treating near-perfect membership as success.

The earlier frozen Python development pilot independently showed the same
qualitative pattern: LR membership $0.63$--$0.88$, split membership $0.99$--$1$,
and split zero exclusion $0$--$0.01$. Its seed family and language RNG differ;
these are separate development runs, not paired trajectories or identical
stochastic replications. Their rates are not pooled. The smaller 32-trajectory
Julia smoke run validates the pipeline and is not the empirical study reported
in Table~\ref{tab:results}.

\subsection{Cross-language and numerical checks}
The recorded Julia suite passes 2,468 assertions, including 2,236 deterministic
scalar comparisons with frozen Python and independently coded base R. These
counts are numerical checks, not independent Monte Carlo replications. The
comparisons cover transition coefficients, conditional means and variances,
likelihood, profiled nuisance estimates, selected LR statistics, and information
calculations. Explicit innovations fix cross-language simulation paths.
Fourteen further assertions in the suite use 256-bit calculations to check
difficult small-KL cases with
relative tolerance $10^{-12}$ and zero absolute tolerance.

The largest absolute discrepancy in closed-form comparisons is
$3.553\times10^{-14}$. Bounded objective/LR comparisons differ by at most
$3.242\times10^{-14}$; bounded parameters and predictive/split scores differ by
at most $4.057\times10^{-6}$, within their predeclared $2\times10^{-5}$ absolute
and relative tolerances. Different scalar optimizers explain why a predictive
score can be more sensitive than the optimized training likelihood. The ten
independent R nuisance optimizations agree in objective to
$6.359\times10^{-13}$. Tolerances were not enlarged after observing failures.

Independent checking found cancellation in the original sampled-KL computation:
forming $u=v_i/\Delta_i-1$ loses its first-order term for tiny
$z=\kappa\Delta_i$. Julia instead evaluates
\[
 u=-z+\frac23z^2-\frac13z^3+\frac{2}{15}z^4+O(z^5)
\]
directly in the small-argument regime and evaluates $u-\log(1+u)$ stably.
For $\kappa=10^{-16}$, $T=1$ and one transition, the old expression gives about
$3.0815\times10^{-33}$, exceeding the path KL of about $2.5\times10^{-33}$.
The corrected calculation gives about $2.5\times10^{-33}$. The frozen references
remain unchanged, and the discrepancy is documented. Even at the fixture value
$\kappa=10^{-10}$, the old cancellation causes about $5.1\times10^{-7}$ relative
error in a KL of order $10^{-20}$; absolute tolerance alone would conceal it.
The separate high-precision checks address that issue. This repair is numerical,
not a new statistical result.

Other tests cover Brownian limits, transition aggregation, affine state and time
rescaling, explicit RNG reproducibility, prefix-only fitting, degeneracy and
failure reporting. Simulator moment checks use 12,000 independent short paths
in each of six cells, with six-standard-error thresholds. Those 72,000 paths
check the simulator; they are not an additional inferential coverage study.

\section{Limitations and next research step}
The empirical design fixes zero drift, zero initial state, unit scale and unit
horizon. It therefore does not establish performance over a nuisance grid,
under longer spans, under endogenous sampling, or with observation noise.
Different $n$ values are not paired on nested paths. One hundred trajectories
per cell permit only limited Monte Carlo precision, and the results are not
locked confirmation evidence.

The study has no faithfully reproduced grid-bootstrap baseline. A comparison
must first specify the paper version, null-imposed nuisance fitting,
initialization, resampling law and inversion statistic~\cite{hansen,lui}.
The known-long-run-mean and $\kappa(\mu-X_t)$ parameterizations in relevant
literature must not be silently identified with unrestricted $a$ at
$\kappa=0$. An ordinary plug-in bootstrap is not automatically finite-sample
exact. The version-specific overlap and remaining full-text access gaps are
recorded in the literature documentation; no novelty conclusion is inferred
from an incomplete search.

Full inversion is also unfinished. The set in \eqref{eq:evalue} need not be a
single finite interval: zero, disconnected components, unbounded tails and
numerical nonresolution must all remain possible. A finite accepted grid is
not a certified continuum set. For any set $C$, half-life uncertainty is
\[
 H(C)=\{\log(2)/\kappa:\kappa\in C,\ \kappa>0\}
 \cup\{+\infty:0\in C\}.
\]
Thus accepting zero retains infinite half-life; a positive unbounded component
approaches half-life zero without attaining it. Numerical tail and crossing
certification is required before reporting finite endpoints or widths.

The next step is a faithful grid-bootstrap reproduction followed by confidence-set
geometry, common-path infill and longer-span comparisons, and a sensitivity
design in $(c,\delta,\text{schedule})$. The current derivations are elementary
specializations of established theory; numerical agreement and exploratory
findings are separate evidence. No new method, efficiency theorem or superiority
claim is established by this report.

\section{Reproducibility}
The repository provides Julia source, a locked project environment, numerical
tests, deterministic reference fixtures and saved development outputs. The report's
tables and plots are generated from those outputs and embedded in this standalone
\LaTeX\ file. It compiles without external figures, data files or bibliography
files. Julia is canonical; Python is a frozen numerical reference and base R
supplies selected independent checks. The public tests use saved synthetic
fixtures and require no Python installation.

\textbf{Run record.} Julia 1.12.5, one Julia thread and one BLAS thread;
100 independent trajectories in each of 16 cells; Xoshiro base seed $2026092803$.
The saved development run is
\begin{center}\small\texttt{@@RUN_ID@@}.\end{center}
Its recorded implementation and environment hashes match the initial source
release \texttt{d08dc35}. The run itself preceded that commit and records its
Git status as uncommitted; the match is established by the saved content hashes.
The 32-path smoke, 1,600-path Python pilot and simulator moment checks have
separate records and purposes.

The README gives installation and execution commands. To regenerate this report's
source from its preserved inputs, run
\begin{center}\small\verb|julia --startup-file=no --project=. scripts/build_report.jl|\end{center}
The report manifest records the template, builder, data and source identities.
New runs create unique output directories; they do not overwrite archived results.
\begin{samepage}
The following SHA-256 fingerprints the declared data and implementation inputs,
not elapsed-time reproducibility or the finished PDF:
\begin{center}\footnotesize\texttt{@@SOURCE_DIGEST@@}\end{center}
\end{samepage}

\begin{thebibliography}{9}
\bibitem{tangchen}
C. Y. Tang and S. X. Chen.
Parameter estimation and bias correction for diffusion processes.
\emph{Journal of Econometrics}, 149(1):65--81, 2009.
\href{https://doi.org/10.1016/j.jeconom.2008.11.001}{doi:10.1016/j.jeconom.2008.11.001}.

\bibitem{yu}
J. Yu.
Bias in the estimation of the mean reversion parameter in continuous time models.
\emph{Journal of Econometrics}, 169(1):114--122, 2012.
\href{https://doi.org/10.1016/j.jeconom.2012.01.004}{doi:10.1016/j.jeconom.2012.01.004}.

\bibitem{hansen}
B. E. Hansen.
The grid bootstrap and the autoregressive model.
\emph{Review of Economics and Statistics}, 81(4):594--607, 1999.
\href{https://doi.org/10.1162/003465399558463}{doi:10.1162/003465399558463}.

\bibitem{lui}
Y. L. Lui, W. Xiao, and J. Yu.
The grid bootstrap for continuous time models.
\emph{Journal of Business \& Economic Statistics}, 40(3):1390--1402, 2022.
\href{https://doi.org/10.1080/07350015.2021.1930014}{doi:10.1080/07350015.2021.1930014}.

\bibitem{universal}
L. Wasserman, A. Ramdas, and S. Balakrishnan.
Universal inference.
\emph{Proceedings of the National Academy of Sciences}, 117(29):16880--16890, 2020.
\href{https://doi.org/10.1073/pnas.1922664117}{doi:10.1073/pnas.1922664117}.
The conditional/profile discussion was checked in
\href{https://arxiv.org/abs/1912.11436v4}{arXiv:1912.11436v4} (2022 revision).
\end{thebibliography}
\end{document}
